% ===== Section 1: "Written all over your face" (language, aphasia, facial expressions) (~16 min) =====
\partaccent{mint}
\topicopener{masks}{Դեմքեր ու խոսք}{Դեմքիդ գրված է}{Ինչու է այդքան դժվար թաքցնել, թե իրականում ինչ ես մտածում}{smug}

% Strangers telling you to smile; faces as mind reading
\begin{frame}{«Մի քիչ ժպտա՛»}
  \begin{columns}[c]
    \begin{column}{0.56\textwidth}
      \begin{itemize}
        \item Նոր ես վիճել կամ ոտքդ դրել ես շան «նվերի» մեջ, ու անծանոթ մեկը քեզ ասում է՝ «Ժպտա՛»
        \item<2-> Դիմախաղը (facial expressions) թույլ է տալիս կարդալ ուրիշի միտքը։ \textbf{Միայն թե դեմքով}
        \item<3-> Ուղեղը ուրիշների հետ շփվելու համար \textbf{շատ} պրոցեսներ ունի
      \end{itemize}
    \end{column}
    \begin{column}{0.4\textwidth}
      \centering
      \begin{tikzpicture}
        \node (m) at (0,0) {\mascot[2.4cm]{sad}};
        \node[say, callout absolute pointer={(m.north east)}, font=\handfont\large] at (1.3,1.8) {Ժպտա՛ մի քիչ};
      \end{tikzpicture}
    \end{column}
  \end{columns}
\end{frame}

% Myth or fact: 90% of communication is non-verbal
\begin{frame}{Առասպել, թե՞ փաստ}
  \mythorfact{Շփման 90\%-ը բառերից դուրս է}{ԿԱԽՎԱԾ Է}{peachdark}{%
    \begin{card}[colback=peach!35!cardwhite]
      Թիվը կախված է նրանից, թե ով է ասում (աղմկոտ ակումբ, թե վանդակ՝ քնած վագրի հետ)։ Բայց շփման \textbf{մեծ մասն} իսկապես բառերից դուրս է։\par
      {\footnotesize «7\% բառ, 38\% ձայն, 55\% դեմք»-ը Մեհրաբյանի փորձերից է, որտեղ մարդիկ \textbf{մեկ բառից} էին կռահում վերաբերմունքը։}
    \end{card}}
  \source{Mehrabian \& Wiener, 1967; Mehrabian, Silent Messages, 1971}
\end{frame}

% The 19th-century map of language: Broca and Wernicke
\begin{frame}{19-րդ դարի խոսքի քարտեզը}
  \begin{columns}[c]
    \begin{column}{0.5\textwidth}
      \centering
      \begin{tikzpicture}
        \node[anchor=south west, inner sep=0] (img) at (0,0) {\includegraphics[width=6.2cm]{language_areas.png}};
        \begin{scope}[x={(img.south east)}, y={(img.north west)}]
          \node[regionlabel, fill=mint!60!cardwhite, anchor=south west] at (0,0.99) {Բրոկայի շրջան};
          \node[regionlabel, fill=lavender!70!cardwhite, anchor=south, visible on=<2->] at (0.81,0.99) {Վերնիկեի շրջան};
          \node[font=\scriptsize, text=plum!70!cream, anchor=north west] at (0.0,0.02) {ձախ կիսագունդ, ճակատը՝ ձախում};
        \end{scope}
      \end{tikzpicture}
    \end{column}
    \begin{column}{0.46\textwidth}
      \begin{card}[colback=mint!35!cardwhite]
        {\blackfont Բրոկայի շրջան} (Broca's area)\par
        Ճակատային բլիթ։ Խոսքի \textbf{կառուցումը}. միտքը ձևակերպել, բառերը ճիշտ դասավորել
      \end{card}
      \onslide<2->{%
      \begin{card}[colback=lavender!40!cardwhite]
        {\blackfont Վերնիկեի շրջան} (Wernicke's area)\par
        Քունքային բլիթ։ Խոսքի \textbf{հասկացումը}. բառերի իմաստը
      \end{card}}
    \end{column}
  \end{columns}
  \vspace{1mm}
  \onslide<3->{{\small Գտել են 19-րդ դարում՝ ուսումնասիրելով մարդկանց, ովքեր վնասվածք էին ստացել \textbf{հենց ճիշտ տեղում}։ Գոնե վնասվածքները իրենք չէին հասցնում (որքան գիտենք)։}\par}
\end{frame}

% Broca's (expressive) aphasia
\begin{frame}{Բրոկայի աֆազիա}
  \begin{columns}[c]
    \begin{column}{0.58\textwidth}
      \begin{itemize}
        \item \textbf{Աֆազիա} (aphasia)՝ խոսքի կամ հասկացողության խանգարում ուղեղի վնասվածքից
        \item Բրոկայի (էքսպրեսիվ) աֆազիա. մարդը \textbf{հասկանում է}, բերանն ու լեզուն առողջ են, բայց սահուն խոսել չի կարողանում
        \item<2-> Հաճախ խանգարված է նաև \textbf{գիրը}։ Բերանն ու ձեռքը տարբեր են, ուրեմն խնդիրը ընդհանուր մասում է՝ \textbf{լեզվի մշակման}
      \end{itemize}
    \end{column}
    \begin{column}{0.38\textwidth}
      \centering
      \begin{tikzpicture}
        \node (m) at (0,0) {\mascot[2.5cm]{thinking}};
        \node[say, callout absolute pointer={(m.north)}, font=\handfont\large] at (0.3,2.2) {Ես... խանութ... հաց...};
      \end{tikzpicture}
    \end{column}
  \end{columns}
\end{frame}

% Wernicke's aphasia: fluent nonsense
\begin{frame}{Վերնիկեի աֆազիա}
  \begin{columns}[c]
    \begin{column}{0.5\textwidth}
      \begin{itemize}
        \item Մարդը \textbf{չի հասկանում} բառերը, թեև զգում է տոնն ու ռիթմը
        \item Խոսում է սահուն ու երկար, բայց իրական ու հորինված բառերի խառնուրդով
        \item<3-> Ու ինքը դա \textbf{չի նկատում}։ Պատկերացրու հիասթափությունը
      \end{itemize}
    \end{column}
    \begin{column}{0.46\textwidth}
      \begin{card}[colback=mint!35!cardwhite]
        {\blackfont\small Ուզում էր ասել}\par
        {\large Գնացի խանութ, հաց գնեցի։}
      \end{card}
      \vspace{2mm}
      \onslide<2->{%
      \begin{card}[colback=lavender!40!cardwhite]
        {\blackfont\small Ասաց}\par
        {\handfont\large Ես գնացիլ խանութ-փանութ այսօրօրօր, գնեցիկ-թնեցիկ հաց-կաց-թաց...}
      \end{card}}
    \end{column}
  \end{columns}
\end{frame}

% Brain scans widened the map
\begin{frame}{Սկաներներից հետո քարտեզը մեծացավ}
  \begin{columns}[c]
    \begin{column}{0.5\textwidth}
      \centering
      \begin{tikzpicture}
        \node[anchor=south west, inner sep=0] (img) at (0,0) {\includegraphics[width=6cm]{language_areas.png}};
        \begin{scope}[x={(img.south east)}, y={(img.north west)}]
          \begin{scope}[visible on=<2->]
            \fill[lavender, opacity=0.75, rotate around={-6:(0.5,0.44)}] (0.5,0.44) ellipse [x radius=1.35cm, y radius=0.18cm];
            \fill[peach, opacity=0.75, rotate around={-8:(0.5,0.32)}] (0.5,0.32) ellipse [x radius=1.35cm, y radius=0.18cm];
            \node[regionlabel, fill=lavender!70!cardwhite] at (0.27,-0.09) {վերին քունքային};
            \node[regionlabel, fill=peach!70!cardwhite] at (0.73,-0.09) {միջին քունքային};
            \node[regionlabel, fill=butter!80!cardwhite] at (0.5,-0.23) {պուտամեն՝ խորքում};
          \end{scope}
          \node[regionlabel, fill=mint!60!cardwhite, anchor=south west] at (0,0.99) {Բրոկա};
          \node[regionlabel, fill=cardwhite, anchor=south] at (0.81,0.99) {Վերնիկե};
        \end{scope}
      \end{tikzpicture}
    \end{column}
    \begin{column}{0.46\textwidth}
      \begin{itemize}
        \item Բրոկան դեռ կարևոր է \textbf{շարահյուսության} (syntax) համար
        \item<2-> Վերնիկեն «պաշտոնազրկվել» է. խոսքը մշակում են քունքային բլթի \textbf{լայն} հատվածներ՝ լսողական կեղևի կողքին (գոնե մի անգամ տրամաբանական է)
        \item<3-> Լեզուն ցրված է ողջ ուղեղով, ուրեմն այն ուղեղի \textbf{հիմնարար} ֆունկցիա է
      \end{itemize}
    \end{column}
  \end{columns}
\end{frame}

% Linguistic relativity and the Sapir-Whorf joke
\begin{frame}{Լեզուն փոխո՞ւմ է մտածողությունը}
  \begin{columns}[T]
    \begin{column}{0.56\textwidth}
      \begin{itemize}
        \item \textbf{Լեզվական հարաբերականություն} (linguistic relativity). լեզուդ որոշում է, թե ինչպես ես մտածում
      \end{itemize}
      \onslide<2->{%
      \begin{bookquote}
        Եթե լեզվում «հուսալի» բառ չլիներ, մարդիկ չէին հասկանա հուսալիությունը ու կաշխատեին \textbf{անշարժ գույքի գործակալ}։
      \end{bookquote}}
      \onslide<3->{{\small Գույների քիչ անուններ ունեցող լեզուների հետազոտությունները \textbf{վիճելի} են։}\par}
    \end{column}
    \begin{column}{0.4\textwidth}
      \onslide<4->{%
      \begin{card}[colback=lavender!40!cardwhite]
        {\blackfont Սեփիր-Ուորֆի վարկած}\par\smallskip
        Սեփիրն ու Ուորֆը միասին ոչինչ չեն գրել ու վարկած չեն առաջարկել։ Վարկածը ծնվեց, երբ հորինեցին նրա \textbf{անունը}. հիանալի օրինակ հենց իր համար։
      \end{card}
      \centering\mascot[1.3cm]{confused}}
    \end{column}
  \end{columns}
\end{frame}

% Neurolinguistic programming
\begin{frame}{Նեյրոլեզվաբանական ծրագրավորում (NLP)}
  \begin{columns}[T]
    \begin{column}{0.5\textwidth}
      \begin{itemize}
        \item Հոգեթերապիայի ու «անձնական աճի» խառնուրդ
        \item Գաղափարը. փոխիր մարդու բառերը, ու կփոխվեն նրա մտածողությունն ու վարքը, ինչպես կոդից «բագ» ջնջելիս \emoji{computer}
        \item<2-> Ապացույցներ գրեթե չկան. \textbf{կեղծ գիտություն} է
      \end{itemize}
    \end{column}
    \begin{column}{0.46\textwidth}
      \onslide<2->{%
      \begin{bookquote}
        Ուղեղն անում է իր ուզածը, ինչ էլ որ աշխարհը շպրտի նրա վրա։ Դժվար թե «շարք կանգնի» մի խնամքով ընտրված արտահայտության առաջ։
      \end{bookquote}}
      \onslide<3->{%
      \begin{card}[colback=mint!35!cardwhite]
        Բայց մի բանում NLP-ն ճիշտ է. \textbf{ոչ բառային} շփումը շատ կարևոր է։
      \end{card}}
    \end{column}
  \end{columns}
\end{frame}

% Oliver Sacks: aphasia patients laughing at the president
\begin{frame}{Նախագահի ելույթը}
  \begin{columns}[c]
    \begin{column}{0.56\textwidth}
      \begin{itemize}
        \item Օլիվեր Սաքս, «Մարդը, ով կնոջը գլխարկի տեղ դրեց» (1985)
        \item Աֆազիայով պացիենտները դիտում են նախագահի ելույթը ու \textbf{քրքջում են}
        \item<2-> Բառերը չհասկանալով՝ նրանք նկատում են այն, ինչ մենք բաց ենք թողնում. դեմքի ջղաձգումներ, ժեստեր, ռիթմ։ \textbf{Ստի նշաններ}
      \end{itemize}
    \end{column}
    \begin{column}{0.38\textwidth}
      \centering
      \bigemoji[1.2cm]{tv}\ \bigemoji[1.2cm]{laughing}\par\medskip
      \onslide<3->{%
      \begin{bookquote}
        Երբ դա գալիս է աշխարհի ամենահզոր մարդուց, մնում է կամ ծիծաղել, կամ լացել։
      \end{bookquote}}
    \end{column}
  \end{columns}
\end{frame}

% Same words, different faces
\begin{frame}{Նույն բառերը, տարբեր դեմքեր}
  \centering
  \begin{tikzpicture}
    \foreach \i/\e/\t/\m/\s in {0/grin/{ուրախ}/{«Շնորհակալ եմ»}/1, 1/rage/{զայրացած}/{«Ինչ ես արել»}/2, 2/nauseated/{զզված}/{«Էս ի՞նչ է»}/3}{
      \begin{scope}[visible on=<\s->]
        \node (f\i) at (\i*4.6,0) {\bigemoji[1.3cm]{\e}};
        \node[say, callout absolute pointer={(f\i.north)}, font=\handfont\large] at (\i*4.6,1.35) {Պետք չէր...};
        \node[font=\blackfont\small, text=plum] at (\i*4.6,-1.0) {\t};
        \node[font=\small, text=plum!75!cream] at (\i*4.6,-1.4) {\m};
      \end{scope}}
  \end{tikzpicture}\par\vspace{1mm}
  \raggedright
  \onslide<4->{%
  \begin{tryit}
    Ասա «պետք չէր» երեք տարբեր դեմքով։ Բառերը նույնն են, իմաստը՝ բոլորովին տարբեր։
  \end{tryit}}
\end{frame}

% Universal expressions, and what we know today
\begin{frame}{Դեմքերը նույնն են ամենուր... թե՞ ոչ}
  \begin{columns}[T]
    \begin{column}{0.5\textwidth}
      \begin{itemize}
        \item Դեմքի արտահայտությունների նկարները ցույց են տվել տարբեր մշակույթների, նաև շատ հեռավոր ժողովուրդների
        \item Գրեթե բոլորը \textbf{ճանաչել են} դրանք
        \item Բըրնեթի եզրակացությունը. դիմախաղը \textbf{բնածին} է։ Ամազոնում ու Նյու Յորքում զարմանում են նույն դեմքով
      \end{itemize}
    \end{column}
    \begin{column}{0.46\textwidth}
      \onslide<2->{%
      \begin{scienceupdate}
        Այսօր սա \textbf{վիճելի} է։ Հին փորձերում ընտրում էին պատրաստի բառերից։ Երբ Նամիբիայի հիմբա ժողովրդին խնդրեցին նկարներն \textbf{ազատ} դասավորել, «համընդհանուր» վեց խումբը չստացվեց։ Դեմքով զգացմունք արտահայտելը շատ է տարբերվում մշակույթից ու իրավիճակից։
      \end{scienceupdate}}
    \end{column}
  \end{columns}
  \source{Gendron et al., Emotion, 2014; Barrett et al., Psychol. Sci. Public Interest, 2019}
\end{frame}

% Emoticons: lines and dots are enough
\begin{frame}{Երկու կետ ու փակագիծ}
  \centering
  \begin{tikzpicture}
    \foreach \i/\s/\t in {0/{:-)}/ուրախություն, 1/{:-(}/տխրություն, 2/{>:-(}/զայրույթ, 3/{:-O}/զարմանք}{
      \node[draw=plum, line width=0.8pt, fill=cardwhite, rounded corners=3mm, minimum width=2.6cm, minimum height=1.4cm,
            font=\blackfont\fontsize{26}{28}\selectfont, text=plum] (e\i) at (\i*3.3,0) {\s};
      \node[font=\small, text=plum, below=1mm of e\i, visible on=<2->] {\t};}
  \end{tikzpicture}\par\vspace{3mm}
  \raggedright
  \begin{columns}[c]
    \begin{column}{0.54\textwidth}
      Պարզ գծեր ու կետեր են, \textbf{նույնիսկ ուղղահայաց չեն}, բայց ուղեղը միանգամից դեմք է տեսնում։ Տեսողական կեղևն ունի դեմքերի համար առանձին հատվածներ (5-րդ գլուխ)։
    \end{column}
    \begin{column}{0.42\textwidth}
      \onslide<3->{%
      \begin{chatbox}
        Գրիր որևէ զգացմունք \textbf{միայն կետադրական նշաններով}։ Մյուսները թող գուշակեն։
      \end{chatbox}}
    \end{column}
  \end{columns}
\end{frame}

% Why faces matter: alarm and feedback
\begin{frame}{Դեմքը՝ ահազանգ ու ուղեցույց}
  \begin{columns}[c]
    \begin{column}{0.54\textwidth}
      \begin{tikzpicture}
        \node[say, text width=6cm, align=left, font=\handfont\normalsize] (s) at (0,0)
          {Չեմ ուզում քեզ անհանգստացնել, բայց կատաղած բորենիների ոհմակ է գալիս մեր վրա...};
        \node[anchor=north, text width=6.4cm, font=\small, text=plum, visible on=<2->] at ([yshift=-3mm]s.south)
          {...մինչև նախադասության վերջը նրանք կհասնեին։ Վախեցած դեմքն \textbf{ավելի արագ} է։};
      \end{tikzpicture}
      \begin{itemize}
        \item<3-> Բոլորը վախեցած են. ուղեղը պատրաստվում է փախչել
        \item<3-> Բոլորը ցնցված են ու զզված. \textbf{անմիջապես կանգնիր}
      \end{itemize}
    \end{column}
    \begin{column}{0.42\textwidth}
      \onslide<4->{%
      \begin{card}[colback=mint!35!cardwhite]
        \emoji{eyes}\ Դեմքեր կարդալիս ակտիվ է \textbf{նշաձև մարմինը} (amygdala)։ Այն մշակում է մեր զգացմունքները ու, երևում է, պետք է նաև ուրիշներինը ճանաչելու համար։ Զզվանքի համար՝ նաև \textbf{պուտամենը}։
      \end{card}}
    \end{column}
  \end{columns}
\end{frame}

% Poker face, until the meteorite
\begin{frame}{Պոկերի դեմք}
  \begin{columns}[c]
    \begin{column}{0.54\textwidth}
      \begin{itemize}
        \item Պոկերիստները պահում են չեզոք (կամ սուտ) դեմք, որ չմատնեն իրենց խաղաքարտերը
        \item 52 խաղաքարտով տարբերակները սահմանափակ են. կարելի է պատրաստվել բոլորին, նույնիսկ անպարտելի «սթրեյթ ֆլեշին»
        \item Երբ գիտես՝ ինչ կարող է լինել, \textbf{գիտակցված} կառավարումը հաղթում է
      \end{itemize}
    \end{column}
    \begin{column}{0.4\textwidth}
      \centering
      \bigemoji[1.2cm]{joker}\par\medskip
      \onslide<2->{%
      \begin{card}[colback=peach!35!cardwhite]
        Բայց եթե խաղի ժամանակ երկնաքարը ծակի առաստաղն ու ընկնի սեղանին, դժվար թե որևէ մեկը կարողանա \textbf{չզարմանալ} \emoji{exploding}
      \end{card}}
    \end{column}
  \end{columns}
\end{frame}

% Voluntary vs involuntary expressions, and the meme
\begin{frame}{Ժպտալ՝ ատամները սեղմած}
  \begin{columns}[c]
    \begin{column}{0.52\textwidth}
      \begin{card}[colback=mint!35!cardwhite]
        {\blackfont Կամային} դիմախաղ. \textbf{շարժիչ կեղև} (motor cortex)։ Ոգևորված դեմք ընկերոջդ ձանձրալի լուսանկարներին նայելիս
      \end{card}
      \begin{card}[colback=peach!35!cardwhite]
        {\blackfont Ակամա} դիմախաղ. \textbf{լիմբիկ համակարգ}, իրական զգացմունքներ
      \end{card}
      \onslide<2->{{\small Նոր կեղևը կարող է ստել, իսկ լիմբիկ համակարգը \textbf{միշտ ազնիվ է}։ Ու նրանք հաճախ են բախվում։}\par}
    \end{column}
    \begin{column}{0.44\textwidth}
      \centering
      \onslide<3->{%
      \begin{meme}[5.4cm]{they_dont_know.png}
        \memecap[text width=2.5cm, font=\blackfont\footnotesize]{0.76}{0.88}{Նրանք չգիտեն, որ ժպիտս կեղծ է}
      \end{meme}}
    \end{column}
  \end{columns}
\end{frame}
